\begin{summarybox}
	\textbf{\textcolor{CSummary}{一句话核心}}

	\textbf{存在取等点时，利用该点导数符号求得参数必要条件，再验证充分性得解。}

	\medskip
	\textbf{\textcolor{CSummary}{使用条件}}
	\begin{itemize}[leftmargin=*, itemsep=0.5em, topsep=0.4em]
		\item \textbf{条件 1（存在取等点）：} 能找到一个 $x_0\in D$ 使 $f(x_0,a)=0$。
		\item \textbf{条件 2（函数可导）：} $f(x,a)$ 在 $D$ 上（至少邻域内）可导。
		\item \textbf{条件 3（不等式含等）：} 不等式为 $\ge0$ 或 $\le0$（若严格不等，需考虑极限情形）。
	\end{itemize}

	\medskip
	\textbf{\textcolor{CSummary}{关键提醒}}
	\begin{itemize}[leftmargin=*, itemsep=0.5em, topsep=0.4em]
		\item \textbf{易错点（端点方向）：} 左端点导数正，右端点导数负（针对 $\ge0$）；反向不等式则符号相反。
		\item \textbf{检查项（必须验证充分性）：} 得出参数范围后要代入验证恒成立，避免失分。
	\end{itemize}
\end{summarybox}
